\documentclass{beamer}
\usepackage{graphicx}
\usepackage{tikz}
\usetikzlibrary{positioning}
\usepackage{amsmath}
\usepackage{amssymb}
\usepackage{amsthm}
\usepackage[table]{xcolor}
\usepackage{caption}
\usepackage{subcaption}
\usepackage{mathastext}

\usetheme{PaloAlto}
\usecolortheme{default}
\usefonttheme{professionalfonts}

\title{Quantile Regression}
\subtitle{Mathematical Formulation and Real World Uses}
\author{Giacomo Boschi}

\date{September 9th, 2026}

\begin{document}
% Title
\begin{frame}
	\titlepage
\end{frame}

% Toc
\begin{frame}
	\frametitle{Table of Contents}
	\tableofcontents
\end{frame}

\section{The Concept}
\begin{frame}{The Foundation}
	\begin{center}
		In OLS, the parameters $(\beta_0,\beta_1)$ are obtained by minimizing the sum of squared residuals, which estimates the conditional mean $E(Y|X)$:
		\[\underset{\beta_0,\beta_1}{min}\ \sum_{i=1}^n(y_i - (\beta_0+\beta_1 x_i))^2\]\hfill \break
		If we replace the squared loss with the absolute value, we get the LAD, which estimates the conditional median $Q_{0,5}(Y|X)$:\hfill \break
		\[\underset{\beta_0,\beta_1}{min} \sum_{i=1}^n\lvert y_i - (\beta_0+\beta_1 x_i)\rvert\]
	\end{center}
\end{frame}

\begin{frame}{From Median to Quantiles through optimization}
	\begin{center}
		The conditional median is a specific quantile; can the loss function be modified to target other quantiles?\hfill \break

		The absolute value is symmetric: it penalises positive and negative residuals equally, this is why we get the median. Therefore to target other quantiles, we have to give different weights to positive and negative residuals:

		\[
			\rho_{\tau}(\epsilon)  = \begin{cases}
				\tau\epsilon\quad     & if\ \epsilon \geq 0 \\
				(\tau-1)\epsilon\quad & if\ \epsilon < 0
			\end{cases}\quad with\ \tau \in (0,1)
		\]
		where $\tau$ is the targeted quantile.
	\end{center}
\end{frame}

\section{Quantile Regression}
\begin{frame}{The Optimization Problem}
	\begin{center}
		Replacing the absolute value with $\rho_\tau$ in the regression problem, we obtain the quantile regression, a generalisation of the LAD, which estimates the conditional quantile $Q_\tau(Y|X)$

		\[
			min_{\beta_0,\beta_1}\ \sum_{i=1}^n \rho_\tau(y_i-(\beta_0+\beta_1x_i))
		\]

		Just like the LAD, this function is piecewise linear and can't be differentiable at 0; one of the approaches is to reformulate and solve it as an LP problem
	\end{center}
\end{frame}

\begin{frame}{The LP Reformulation: splitting the residuals}
	\begin{center}
		Each residual $\epsilon = y_i - (\beta_0+\beta_1x_i)$ can be decomposed into its positive and negative parts:

		\[
			\epsilon_i=u_i-v_i,\quad with\ u_i, v_i \geq 0
		\]
		\[
			\begin{cases}
				\epsilon_i \geq 0\colon &
				\begin{aligned}[t]
					 & u_i = \epsilon_i,\; v_i = 0       \\
					 & \Rightarrow \rho_\tau(\epsilon_i)
					= \tau\epsilon_i
					= \tau u_i + (1-\tau)\cdot 0
				\end{aligned} \\
				\epsilon_i < 0\colon    &
				\begin{aligned}[t]
					 & u_i = 0,\; v_i = -\epsilon_i      \\
					 & \Rightarrow \rho_\tau(\epsilon_i)
					= (\tau-1)\epsilon_i
					= \tau\cdot 0 + (1-\tau)v_i
				\end{aligned}
			\end{cases}
		\]

		In both cases: $\rho_\tau(\epsilon_i) = \tau u_i + (1-\tau)v_i$
	\end{center}
\end{frame}
\end{document}
